% =============================================================================
% 061 / 068
% Status: raw
% =============================================================================
% Notes:
%
% Source question:
% 當代-吳越民族，
% 是否要面對 五代宋-江淮類似的命運？
% (作為東亞大陸社會相對健全的地區，
% 即將面對經濟文化不可逆的重傷？)
% 
% 江淮盛衰期：
% 南北朝(-傷)
% 隋唐(++++極盛)
% 五代(-衰)
% 宋(- -再衰)
% 明(- - -大洪水)
% 明以後到清末
% (- 自然與社會生態被京杭運河嚴重破壞) 
% 清末
% (- -運河失效) 
% 
% 吳越盛衰期：
% 隋唐(+小盛)-
% 五代(+++極盛)
% 宋元(++小盛)
% 明初(- -小洪水)
% 晚明 清初(+小盛)
% 太平天國(- - 小洪水)
% 晚清民初(+小盛)
% 抗戰(- 衰)-
% CCP初(- -小洪水)
% 改革開放-2018年
% (++經濟大盛,文化衰)  
% 中國衰退,臺海戰爭前線
% (- - -潛在經濟、文化大衰退的可能？)
% =============================================================================

\chapter[當代-吳越民族， 是否要面對 五代宋-江淮類似的命運？ (作為東亞大陸社會相對健全的地區， 即將面對經濟文化不可逆的重傷？) 江淮盛衰期： 南北朝(-傷) 隋唐(++++極盛) 五代(-衰) 宋(- -再衰) 明(- - -大洪水) 明以後到清末 (- 自然與社會生態被京杭運河嚴重破壞) 清末 (- -運河失效) 吳越盛衰期： 隋唐(+小盛)- 五代(+++極盛) 宋元(++小盛) 明初(- -小洪水) 晚明 清初(+小盛) 太平天國(- - 小洪水) 晚清民初(+小盛) 抗戰(- 衰)- CCP初(- -小洪水) 改革開放-2018年 (++經濟大盛,文化衰) 中國衰退,臺海戰爭前線 (- - -潛在經濟、文化大衰退的可能？)]{當代-吳越民族， \\[0.35em] 是否要面對 五代宋-江淮類似的命運？ \\[0.35em] (作為東亞大陸社會相對健全的地區， \\[0.35em] 即將面對經濟文化不可逆的重傷？) \\[0.35em] 江淮盛衰期： \\[0.35em] 南北朝(-傷) \\[0.35em] 隋唐(++++極盛) \\[0.35em] 五代(-衰) \\[0.35em] 宋(- -再衰) \\[0.35em] 明(- - -大洪水) \\[0.35em] 明以後到清末 \\[0.35em] (- 自然與社會生態被京杭運河嚴重破壞) \\[0.35em] 清末 \\[0.35em] (- -運河失效) \\[0.35em] 吳越盛衰期： \\[0.35em] 隋唐(+小盛)- \\[0.35em] 五代(+++極盛) \\[0.35em] 宋元(++小盛) \\[0.35em] 明初(- -小洪水) \\[0.35em] 晚明 清初(+小盛) \\[0.35em] 太平天國(- - 小洪水) \\[0.35em] 晚清民初(+小盛) \\[0.35em] 抗戰(- 衰)- \\[0.35em] CCP初(- -小洪水) \\[0.35em] 改革開放-2018年 \\[0.35em] (++經濟大盛,文化衰) \\[0.35em] 中國衰退,臺海戰爭前線 \\[0.35em] (- - -潛在經濟、文化大衰退的可能？)}
\qaanswer{
如果保護不了自己的勝利果實，那之前的繁盛榮昌都是虛假的。
}

